\documentclass[11pt,a4paper]{article}

% Required packages
\usepackage[utf8]{inputenc}
\usepackage[T1]{fontenc}
\usepackage{amsmath}
\usepackage{graphicx}
\usepackage{booktabs}
\usepackage{hyperref}
\usepackage[margin=1in]{geometry}
\usepackage{natbib}
\usepackage{xcolor}
\usepackage{lineno}
\usepackage{authblk}

\bibliographystyle{unsrtnat}

\title{The SUMO E3 Ligase Ema2 Promotes Genome-Wide lncRNA Transcription from the Somatic Macronucleus to Enable scnRNA-Directed DNA Elimination in \textit{Tetrahymena}}

\author[1]{Author A}
\author[1]{Author B}
\author[1,2]{Author C}
\affil[1]{Department of Molecular Biology, University of Rochester, Rochester, NY, USA}
\affil[2]{Corresponding author: \texttt{author.c@rochester.edu}}

\date{}

\begin{document}
\maketitle

\begin{abstract}
Small RNA-directed chromatin silencing pathways require nascent long non-coding RNA (lncRNA) transcription from target loci, yet the molecular mechanisms that actively promote this transcription remain poorly understood. In \textit{Tetrahymena thermophila}, programmed DNA elimination removes approximately 12,000 internal eliminated sequences (IESs) from the developing somatic macronucleus (MAC), guided by small RNAs (scnRNAs) produced from the germline micronucleus (MIC). Here we identify Ema2, a conjugation-specific SUMO E3 ligase containing an SP-RING domain, as an essential factor promoting lncRNA transcription from the parental MAC during mid-conjugation. Ema2 interacts directly with the SUMO E2 conjugating enzyme Ubc9 and promotes SUMOylation of the transcription regulator Spt6. Knockout of \textit{EMA2} abolishes lncRNA accumulation from the parental MAC, impairs target-directed scnRNA degradation (TDSD), disrupts heterochromatin formation in the new MAC, and partially blocks DNA elimination. Ema2 is required for the chromatin association of both Spt6 and RNA Polymerase II (RNAPII), suggesting a mechanism by which SUMOylation maintains transcription elongation competence. However, SUMOylation-defective Spt6 mutants do not phenocopy the lncRNA defect, indicating that Ema2 has additional SUMOylation targets critical for lncRNA transcription. Our findings reveal a previously unrecognized role for SUMOylation in connecting lncRNA production to small RNA-mediated genome surveillance and establish a molecular framework for understanding how active transcription from regions destined for silencing is temporally controlled.
\end{abstract}

\section{Introduction}

Small RNA-directed chromatin silencing is a conserved mechanism through which organisms defend genome integrity against transposable elements and other repetitive sequences \citep{Malone2009, Aravin2007, Castel2013}. A fundamental paradox in this process is that the very regions targeted for silencing must be actively transcribed to produce the long non-coding RNA (lncRNA) scaffolds required for small RNA effector recruitment \citep{Grewal2007, Moazed2009}. In most model organisms, including fission yeast, \textit{Drosophila}, and plants, the source and target loci of small RNAs overlap, making it difficult to dissect the mechanisms that promote lncRNA transcription independently of its role in small RNA biogenesis \citep{Volpe2002, Brennecke2007, Matzke2015}.

\textit{Tetrahymena thermophila} provides a uniquely powerful system for studying this process because it maintains two functionally distinct nuclei within the same cell: a germline micronucleus (MIC) and a somatic macronucleus (MAC) \citep{Karrer2012, Chalker2011}. During conjugation, the sexual reproduction phase of the \textit{Tetrahymena} life cycle, the MIC undergoes meiosis and approximately 12,000 internal eliminated sequences (IESs) are removed from the developing new MAC \citep{Mochizuki2010, Yao2003}. This programmed DNA elimination reduces the genome from approximately 200~Mb (MIC) to approximately 103~Mb (MAC) and is guided by scan RNAs (scnRNAs), a class of approximately 26--31~nt small RNAs produced from the MIC genome during meiosis \citep{Mochizuki2002, Liu2007}.

The current model for DNA elimination proceeds through a multi-step scanning mechanism \citep{Mochizuki2002, Aronica2008, Noto2015}. First, scnRNAs are generated from the MIC and loaded onto the Piwi protein Twi1 during early conjugation. These Twi1-scnRNA complexes then translocate to the parental MAC, where they scan lncRNAs transcribed from the somatic genome. scnRNAs complementary to MAC-destined sequences (MDS) are selectively degraded through a process termed target-directed scnRNA degradation (TDSD), while IES-matching scnRNAs are retained \citep{Mochizuki2004, Noto2015}. The remaining IES-specific scnRNAs are then transferred to the developing new MAC, where they guide deposition of heterochromatin marks (H3K27me3 and H3K9me3) and ultimately direct DNA elimination \citep{Liu2007, Taverna2013}.

A critical step in this pathway is the transcription of lncRNAs from the parental MAC during mid-conjugation. These lncRNAs serve as the scanning template against which scnRNAs are compared. Despite the centrality of this step, the molecular mechanism that promotes genome-wide lncRNA transcription from the parental MAC had remained unknown. SUMOylation, the covalent attachment of the Small Ubiquitin-like Modifier (SUMO) protein to target lysines, is a versatile post-translational modification implicated in transcription regulation, chromatin remodeling, and DNA repair \citep{Gareau2010, Flotho2013, Raman2013}. However, its role in lncRNA transcription for genome surveillance had not been established.

Here we identify Ema2, a conjugation-specific protein containing an SP-RING domain characteristic of SUMO E3 ligases, as an essential factor for lncRNA transcription from the parental MAC. Our findings connect the SUMOylation machinery to the small RNA-directed DNA elimination pathway and reveal how transcription from regions destined for silencing is temporally controlled during conjugation.

\section{Results}

\subsection{Ema2 is a conjugation-specific protein that localizes to the macronucleus}

To identify factors involved in DNA elimination, we searched for conjugation-specific genes with predicted roles in chromatin regulation. Analysis of microarray expression data \citep{Miao2009} revealed that \textit{EMA2} (TTHERM\_00113330) is undetectable in growing and starved cells but is highly expressed during conjugation, peaking between 3 and 6 hours post-mixing (hpm) (Table~\ref{tab:expression}).

\begin{table}[htbp]
\centering
\caption{EMA2 mRNA expression profile across growth and developmental stages.}
\label{tab:expression}
\begin{tabular}{ll}
\toprule
\textbf{Growth/Stage Condition} & \textbf{EMA2 mRNA Expression} \\
\midrule
Growing cells (low density) & Undetectable \\
Growing cells (medium density) & Undetectable \\
Growing cells (high density) & Undetectable \\
Starved cells (0--24 hr) & Undetectable \\
Conjugation 0 hpm & Begins to appear \\
Conjugation 3--6 hpm & High expression \\
Conjugation 8--12 hpm & Present \\
Conjugation 14--18 hpm & Declining \\
\bottomrule
\end{tabular}
\end{table}

To determine the subcellular localization of Ema2, we generated strains expressing HA-tagged Ema2 from the endogenous locus. Immunofluorescence microscopy revealed a dynamic localization pattern (Table~\ref{tab:localization}): Ema2-HA first appeared in the parental MAC at 3~hpm, persisted through 6~hpm, then switched to the developing new MAC at approximately 8~hpm before fading by 14~hpm. This parental-to-new MAC transition is reminiscent of the localization dynamics of the Piwi protein Twi1 and the PRC2 complex \citep{Noto2015, Liu2007}.

\begin{table}[htbp]
\centering
\caption{Ema2-HA protein localization during conjugation.}
\label{tab:localization}
\begin{tabular}{lll}
\toprule
\textbf{Time Point} & \textbf{Ema2-HA Location} & \textbf{Notes} \\
\midrule
Vegetative & Not detectable & No expression in growing cells \\
3 hpm & Parental MAC & First appearance \\
6 hpm & Parental MAC & Continued localization \\
8 hpm & New MAC & Localization switch \\
12 hpm & New MAC (fading) & Signal diminishing \\
14 hpm & New MAC (fading) & Signal further reduced \\
\bottomrule
\end{tabular}
\end{table}

\subsection{Ema2 is required for DNA elimination and produces inviable progeny}

To determine whether Ema2 is required for DNA elimination, we generated somatic knockout strains by disrupting all MAC copies of \textit{EMA2}. DNA elimination was assessed by fluorescence in situ hybridization (FISH) using probes against the Tlr1 transposable element at 36~hpm \citep{Mochizuki2002}. In wild-type cells, Tlr1 signal was detected only in the MIC, indicating successful elimination from the new MAC. In contrast, \textit{EMA2} knockout cells retained Tlr1 signal in the new MAC, although the signal was weaker than that observed in \textit{TWI1} knockout controls, indicating partial rather than complete blocking of DNA elimination (Table~\ref{tab:fish}).

\begin{table}[htbp]
\centering
\caption{FISH analysis of DNA elimination using Tlr1 probes at 36 hpm.}
\label{tab:fish}
\begin{tabular}{llll}
\toprule
\textbf{Genotype} & \textbf{Tlr1 in MIC} & \textbf{Tlr1 in New MAC} & \textbf{Interpretation} \\
\midrule
Wild-type & Detected & Not detected & Elimination complete \\
\textit{EMA2} KO & Detected & Detected (reduced) & Partially blocked \\
\textit{TWI1} KO & Detected & Detected (strong) & Completely blocked \\
\bottomrule
\end{tabular}
\end{table}

Consistent with the DNA elimination defect, \textit{EMA2} knockout strains produced no viable sexual progeny, confirming the essential role of Ema2 in the DNA elimination pathway.

\subsection{Ema2 is required for target-directed scnRNA degradation}

The partial DNA elimination defect in \textit{EMA2} knockout cells suggested a deficiency upstream of heterochromatin deposition. We therefore examined TDSD by northern blot hybridization using a probe complementary to MDS regions (Mi-9). In wild-type crosses, MDS-complementary scnRNAs were progressively degraded between 3 and 6~hpm, consistent with functional TDSD. In \textit{EMA2} knockout crosses, these scnRNAs persisted through 6~hpm (Table~\ref{tab:tdsd}), indicating that TDSD was blocked. This phenotype was rescued by expression of HA-tagged Ema2, confirming specificity. Furthermore, H3K27me3 deposition in the new MAC was reduced or absent in knockout cells, consistent with impaired heterochromatin formation downstream of the TDSD defect \citep{Liu2007}.

\begin{table}[htbp]
\centering
\caption{Target-directed scnRNA degradation (TDSD) assessed by northern blot.}
\label{tab:tdsd}
\begin{tabular}{lllll}
\toprule
\textbf{Genotype} & \textbf{3 hpm} & \textbf{4.5 hpm} & \textbf{6 hpm} & \textbf{Trend} \\
\midrule
WT & Present & Reduced & Strongly reduced & Progressive TDSD \\
\textit{EMA2} KO & Present & Persists & Persists & TDSD blocked \\
\bottomrule
\end{tabular}
\end{table}

\subsection{Ema2 is a SUMO E3 ligase that interacts with Ubc9}

Sequence analysis revealed that Ema2 contains an SP-RING domain with conserved cysteine and histidine residues required for zinc coordination, a hallmark of the PIAS/Siz family of SUMO E3 ligases \citep{Gareau2010, Flotho2013}. Alignment with known SUMO E3 ligases from diverse organisms (Table~\ref{tab:spring}) confirmed the conservation of key residues.

\begin{table}[htbp]
\centering
\caption{SP-RING domain comparison across species.}
\label{tab:spring}
\begin{tabular}{llll}
\toprule
\textbf{Organism} & \textbf{Protein} & \textbf{Domain} & \textbf{Conserved Residues} \\
\midrule
\textit{Tetrahymena} & Ema2 & SP-RING & Conserved Cys/His \\
\textit{S. cerevisiae} & MMS21 & SP-RING & Conserved \\
\textit{S. cerevisiae} & SIZ1 & SP-RING & Conserved \\
\textit{D. melanogaster} & Su(var)2-10 & SP-RING & Conserved \\
\textit{H. sapiens} & PIAS1 & SP-RING & Conserved \\
\bottomrule
\end{tabular}
\end{table}

To test whether Ema2 functions as a SUMO E3 ligase, we performed in vitro binding assays with recombinant proteins. GST-Ema2, but not GST alone, co-purified with His-Ubc9 (the SUMO E2 conjugating enzyme), demonstrating a direct physical interaction. In reconstituted SUMOylation assays, addition of Ema2 to E1 and E2 enzymes enhanced SUMOylation activity, confirming its function as a SUMO E3 ligase.

\subsection{Ema2 promotes SUMOylation of the transcription regulator Spt6}

To identify in vivo SUMOylation targets of Ema2, we expressed His-tagged Smt3 (SUMO) in wild-type and \textit{EMA2} knockout backgrounds and purified His-Smt3-conjugated proteins by Ni-NTA pulldown from conjugating cells at 6~hpm. Label-free quantification (LFQ) mass spectrometry identified Spt6, a conserved transcription regulator, as a protein with significantly reduced SUMOylation in \textit{EMA2} knockout cells (Table~\ref{tab:sumo}).

\begin{table}[htbp]
\centering
\caption{Ema2-dependent SUMOylation targets identified by mass spectrometry.}
\label{tab:sumo}
\begin{tabular}{lll}
\toprule
\textbf{Protein} & \textbf{Ema2-Dependent?} & \textbf{Evidence} \\
\midrule
Spt6 & Yes & Higher LFQ intensity in WT vs. KO \\
Other proteins & Not significant & Possibly masked \\
\bottomrule
\end{tabular}
\end{table}

Western blot analysis of Spt6-HA revealed slower-migrating species appearing at 4.5--6~hpm in wild-type cells, consistent with conjugation-dependent SUMOylation. These species were not observed in vegetative cells. The functionality of the tagged Smt3 constructs was confirmed by rescue experiments: both His-Smt3 and HA-Smt3 restored viability to \textit{SMT3} knockout progeny, whereas empty vector did not \citep{Gareau2010}.

\subsection{Ema2 is required for lncRNA accumulation from the parental MAC}

We next tested whether Ema2 is required for lncRNA production from the parental MAC. Using RT-PCR primers designed to span MDS-IES junctions (which detect lncRNAs that read through IES boundaries but not processed mRNAs), we found that lncRNA products were readily detected in wild-type crosses at 6~hpm but were absent or strongly reduced in \textit{EMA2} knockout crosses (Table~\ref{tab:lncrna}). Genomic DNA PCR controls confirmed that the target loci were intact in knockout strains, ruling out genomic deletion as the cause.

\begin{table}[htbp]
\centering
\caption{lncRNA accumulation assessed by RT-PCR at 6 hpm.}
\label{tab:lncrna}
\begin{tabular}{lll}
\toprule
\textbf{lncRNA Region} & \textbf{WT} & \textbf{\textit{EMA2} KO} \\
\midrule
MDS-IES junction products & Detected & Not detected or strongly reduced \\
\bottomrule
\end{tabular}
\end{table}

Immunofluorescence staining with the J2 antibody (anti-long dsRNA) further confirmed that dsRNA accumulation in the parental MAC was strongly reduced in \textit{EMA2} knockout cells compared to wild-type.

\subsection{Ema2 promotes chromatin association of Spt6 and RNAPII}

To determine how Ema2 affects the transcription machinery, we examined the chromatin association of Spt6 and the RNAPII subunit Rpb3 by immunofluorescence. In wild-type crosses, both Spt6-HA and Rpb3 showed strong chromatin association during mid-conjugation. In \textit{EMA2} knockout crosses, chromatin association of both proteins was substantially reduced, with the majority of Rpb3 lost from chromatin (Table~\ref{tab:chromatin}).

\begin{table}[htbp]
\centering
\caption{Chromatin association of transcription factors in WT and \textit{EMA2} KO.}
\label{tab:chromatin}
\begin{tabular}{lll}
\toprule
\textbf{Cross} & \textbf{Spt6-HA Signal} & \textbf{Rpb3 Signal} \\
\midrule
WT & Strong chromatin association & Strong association \\
\textit{EMA2} KO & Reduced & Majority lost \\
\bottomrule
\end{tabular}
\end{table}

\subsection{SUMOylation-defective Spt6 does not phenocopy the Ema2 knockout}

To test whether Spt6 SUMOylation is the sole mechanism through which Ema2 promotes lncRNA transcription, we generated SUMOylation-defective Spt6 mutants. Surprisingly, these mutants showed normal lncRNA production and normal DNA elimination (Table~\ref{tab:spt6mut}), in contrast to the strong defects observed in \textit{EMA2} knockout cells.

\begin{table}[htbp]
\centering
\caption{Comparison of Spt6 SUMOylation-defective mutants and \textit{EMA2} KO.}
\label{tab:spt6mut}
\begin{tabular}{lll}
\toprule
\textbf{Spt6 Variant} & \textbf{lncRNA Production} & \textbf{DNA Elimination} \\
\midrule
Wild-type Spt6 & Normal & Normal \\
SUMOylation-defective Spt6 & Normal & Normal \\
\textit{EMA2} KO & Impaired & Partially blocked \\
\bottomrule
\end{tabular}
\end{table}

This result indicates that Ema2 has additional SUMOylation targets beyond Spt6 that are critical for lncRNA transcription from the parental MAC.

\section{Discussion}

Our study reveals that the SUMO E3 ligase Ema2 is an essential component of the molecular machinery that promotes genome-wide lncRNA transcription from the parental MAC during conjugation in \textit{Tetrahymena}. This lncRNA transcription is a prerequisite for the scanning mechanism that underlies programmed DNA elimination \citep{Mochizuki2002, Noto2015}. The identification of Ema2 provides the first direct link between SUMOylation and lncRNA transcription in the context of small RNA-mediated genome surveillance.

The pathway model that emerges from our data (Table~\ref{tab:pathway}) positions Ema2-dependent lncRNA transcription as a central step between scnRNA biogenesis and the downstream events of heterochromatin formation and DNA elimination. This temporal ordering is consistent with the conjugation-specific expression and dynamic MAC localization of Ema2.

\begin{table}[htbp]
\centering
\caption{Pathway summary: steps in scnRNA-directed DNA elimination.}
\label{tab:pathway}
\begin{tabular}{llll}
\toprule
\textbf{Step} & \textbf{Event} & \textbf{Location} & \textbf{Timing} \\
\midrule
1 & scnRNA biogenesis from MIC & MIC & Early conjugation \\
2 & scnRNAs loaded onto Twi1 & Cytoplasm & Early conjugation \\
3 & lncRNA transcription (Ema2-dependent) & Parental MAC & Mid-conjugation \\
4 & scnRNAs scan parental MAC lncRNAs & Parental MAC & Mid-conjugation \\
5 & TDSD: MAC-matching scnRNAs degraded & Parental MAC & Mid-conjugation \\
6 & IES-matching scnRNAs transferred & New MAC & 8 hpm onward \\
7 & H3K27me3/H3K9me3 deposition & New MAC & 8--14 hpm \\
8 & Programmed DNA elimination & New MAC & 14--18 hpm \\
\bottomrule
\end{tabular}
\end{table}

The finding that SUMOylation-defective Spt6 mutants do not phenocopy the \textit{EMA2} knockout is particularly informative. It suggests that Ema2 SUMOylates multiple targets to promote lncRNA transcription, and that the loss of Spt6 SUMOylation alone is insufficient to recapitulate the full phenotype. This is consistent with the broad substrate specificity often observed for SUMO E3 ligases \citep{Flotho2013, Raman2013}. Future proteomic studies with deeper coverage may reveal additional Ema2 substrates.

The mechanism by which Ema2-dependent SUMOylation promotes RNAPII chromatin association remains to be fully elucidated. SUMOylation of transcription factors can modulate their stability, localization, and protein-protein interactions \citep{Gareau2010}. Given that Spt6 functions as a histone chaperone that facilitates transcription elongation by maintaining chromatin structure \citep{Bortvin1996, Ivanovska2011}, SUMOylation of Spt6 may enhance its ability to remodel chromatin for RNAPII passage. However, since Spt6 SUMOylation alone is dispensable, additional mechanisms must be invoked.

Our findings also highlight the utility of \textit{Tetrahymena} as a model for studying the interplay between lncRNA transcription and small RNA pathways. The spatial separation of source (MIC) and target (MAC) nuclei provides an unparalleled opportunity to dissect molecular mechanisms that are difficult to study in other organisms where these compartments overlap \citep{Chalker2011, Karrer2012}.

\section{Materials and Methods}

\subsection{Strains and culture conditions}
\textit{Tetrahymena thermophila} wild-type strains B2086, CU427, and CU428 were used. MIC-defective star strains B*VI and B*VII were used where indicated. Cells were grown in SPP medium containing 2\% proteose peptone at 30$^\circ$C. For conjugation, growing cells ($\sim$5--7 $\times$ 10$^5$/mL) were washed and pre-starved for 12--24 hours, then mixed in 10~mM Tris pH 7.5 at 30$^\circ$C \citep{Miao2009}.

\subsection{Gene knockout and tagging}
Somatic knockout of \textit{EMA2} was achieved by disrupting all MAC copies. For localization studies, the pEMA2-HA-neo3 construct was used for C-terminal HA tagging at the endogenous locus, selected with 100~$\mu$g/mL paromomycin and 1~$\mu$g/mL CdCl$_2$.

\subsection{Antibodies}
The following antibodies were used: anti-H3K27me3 (Millipore 07-449), anti-HA HA11 (Covance clone 6B12), anti-H3K9me3 (Active Motif 39162), anti-GST (BD Biosciences 554805), anti-His (Proteintech 66005-1-Ig), anti-long dsRNA J2 (Jena Bioscience), anti-$\alpha$-tubulin 12G10 (DSHB), anti-Rpb3, and anti-Smt3 (Table~\ref{tab:antibodies}).

\begin{table}[htbp]
\centering
\caption{Antibodies used in this study.}
\label{tab:antibodies}
\small
\begin{tabular}{lllll}
\toprule
\textbf{Antibody} & \textbf{Target} & \textbf{Source} & \textbf{IF} & \textbf{WB} \\
\midrule
Anti-H3K27me3 & H3K27me3 & Millipore 07-449 & 1:1000 & 1:10000 \\
Anti-HA HA11 & HA tag & Covance 6B12 & 1:1000 & 1:10000 \\
Anti-H3K9me3 & H3K9me3 & Active Motif 39162 & 1:1000 & 1:10000 \\
Anti-GST & GST tag & BD 554805 & -- & 1:10000 \\
Anti-His & His tag & Proteintech & -- & 1:10000 \\
Anti-dsRNA J2 & Long dsRNA & Jena Bioscience & 1:1000 & -- \\
Anti-$\alpha$-tubulin & Tubulin & DSHB 12G10 & -- & 1:10000 \\
Anti-Rpb3 & RNAPII & Previously described & 1:1000 & 1:10000 \\
Anti-Smt3 & SUMO & Gift & -- & 1:10000 \\
\bottomrule
\end{tabular}
\end{table}

\subsection{Fluorescence in situ hybridization (FISH)}
Cells were fixed at 36~hpm and hybridized with fluorescently labeled probes against the Tlr1 transposable element \citep{Mochizuki2002}.

\subsection{Northern blot analysis}
Total RNA was extracted at 3, 4.5, and 6~hpm. Small RNA fractions were resolved on denaturing polyacrylamide gels and probed with radiolabeled oligonucleotides complementary to MDS regions (Mi-9 probe).

\subsection{In vitro interaction and SUMOylation assays}
Recombinant GST-Ema2 and His-Ubc9 were expressed in \textit{E. coli} and purified. GST pulldown assays were performed to assess direct interaction. For in vitro SUMOylation, E1, E2 (Ubc9), SUMO (Smt3), and ATP were incubated with or without Ema2, and products were analyzed by western blot.

\subsection{Mass spectrometry}
His-Smt3 conjugates were purified by Ni-NTA pulldown from conjugating cells at 6~hpm in wild-type and \textit{EMA2} knockout backgrounds. Proteins were identified by LC-MS/MS with label-free quantification \citep{Cox2014}.

\subsection{RT-PCR}
Primers spanning MDS-IES junctions were designed to detect lncRNAs that read through IES boundaries. RT-PCR was performed on total RNA from 6~hpm conjugating cells. Genomic DNA PCR with the same primers served as a positive control.

\begin{table}[htbp]
\centering
\caption{Key molecular players and functions in the DNA elimination pathway.}
\label{tab:players}
\begin{tabular}{lll}
\toprule
\textbf{Protein} & \textbf{Function} & \textbf{Role} \\
\midrule
Ema2 & SUMO E3 ligase & Promotes lncRNA transcription \\
Ubc9 & SUMO E2 enzyme & Interaction partner of Ema2 \\
Smt3 & SUMO protein & Conjugated to targets \\
Spt6 & Transcription regulator & SUMOylated by Ema2 \\
Rpb3 & RNAPII subunit & Required for lncRNA transcription \\
Twi1 & Piwi protein & Mediates TDSD \\
H3K27me3 & Histone modification & Marks IESs for elimination \\
H3K9me3 & Histone modification & Heterochromatin mark \\
\bottomrule
\end{tabular}
\end{table}

\begin{table}[htbp]
\centering
\caption{\textit{Tetrahymena} genome parameters.}
\label{tab:genome}
\begin{tabular}{ll}
\toprule
\textbf{Parameter} & \textbf{Value} \\
\midrule
MIC genome size & $\sim$200 Mb \\
MAC genome size & $\sim$103 Mb \\
Number of IESs eliminated & $\sim$12,000 \\
DNA eliminated & $\sim$97 Mb \\
scnRNA size & $\sim$26--31 nt \\
Normal elimination completion & $\sim$14--18 hpm \\
\bottomrule
\end{tabular}
\end{table}

\section*{Acknowledgments}
We thank members of the laboratory for helpful discussions and for reagents including the anti-Smt3 antibody.

\bibliography{references}

\end{document}
